% Propietario conservado: /Users/ruben/Documents/New project/output/REVISION_ARGUMENTAL_APERTURAS_CIERRES_20260915/VI/delivery/MANUSCRIPT_VI_ES_EN_COMPLETE/source_es/sections/03_conjugacion_mobius.tex
\section{Conjugation, orientation, and return of material sections}
\label{vi:vi:sec:conjugacion}

The persistent identity of Section~\ref{vi:vi:sec:particula} admits
inversion operations at different levels of its structure.
Cellular inversion acts on the atlas; dodecaphase conjugation
acts on the orientation record; lifting a loop acts
on a fibre; a spinor rotation acts on a representation
space. Specifying these domains allows the operations to be
composed and explains what information is preserved in each return.

\subsection{Dodecaphase involution and parity of observables}
\label{vi:vi:subsec:conjugacion-registro}

Let $\mathcal T_{12}=\{-1,0,1\}^{12}$ be the ambient space of tritic
records with twelve positions. The admissible TPK domain is the
subset satisfying the route compatibilities already constructed. The
operation on the ambient space is
\begin{equation}
 C_M:\mathcal T_{12}\longrightarrow\mathcal T_{12},\qquad
 (C_M\tau)_j=-\tau_{11-j},\quad 0\leq j<12.
 \label{vi:vi:eq:cm}
\end{equation}
Restriction to admissible records uses stability of that
domain under conjugation. Reversed order and sign reversal
are simultaneous components of the operation.

\begin{proposition}[Linear and quadratic parity of the record]
\label{vi:vi:prop:paridad-registro}
We have $C_M^2=I$. If $a_j=a_{11-j}$ and
$L_a(\tau)=\sum_ja_j\tau_j$, then
$L_a(C_M\tau)=-L_a(\tau)$. If a quadratic form
$Q_b(\tau)=\sum_{j,k}b_{jk}\tau_j\tau_k$ satisfies
$b_{jk}=b_{11-j,11-k}$, then $Q_b(C_M\tau)=Q_b(\tau)$.
\end{proposition}
\begin{proof}
At each position,
$[C_M(C_M\tau)]_j=-(-\tau_{11-(11-j)})=\tau_j$.
For the linear functional, substituting $k=11-j$ in the sum gives
$L_a(C_M\tau)=-\sum_ka_{11-k}\tau_k=-L_a(\tau)$.
In the quadratic expression the two sign reversals cancel, and
the substitution $(r,s)=(11-j,11-k)$ recovers $Q_b$ through the
symmetry of its coefficients.
\end{proof}

Linear orientation and memory correlations therefore have
different laws. Even a composition of these observables must
preserve their parities: if $L_a$ is odd, its exponential character
satisfies
\begin{equation}
 \exp(L_a(C_M\tau))=\exp(L_a(\tau))^{-1}.
 \label{vi:vi:eq:caracter-inverso}
\end{equation}
The equality expresses multiplicative inversion. Invariance of a
complete material operator is studied on that operator, with its
channels and representation, and is established by the corresponding
covariance.

Cellular inversion in Theorem~\ref{vi:vi:thm:selector-central} is
$\rho_c:\mathcal C_{81}\to\mathcal C_{81}$, whereas $C_M$ has
domain \eqref{vi:vi:eq:cm}. Both are involutions, but comparing
them requires the route record map and its transport
law. This datum prevents conjugation of a history from being replaced
by inversion of an isolated atlas coordinate.

\subsection{Lifting to antisections and preservation of mass}
\label{vi:vi:subsec:antisecciones}

Conjugation is applied to the particle through two compatible
maps. The base map $C_m^{\mathrm b}$ transports the state
$x_m$ to its conjugate state $\bar x_m$; the lift
$\widetilde C_m:\mathscr E_m(x)\to\mathscr E_m(\bar x)$
transports its material sections. Let $t_m$ be the reader of the
dodecaphase register and $\bar p_{n,m}$ the restrictions in the
conjugate family. Their conditions are preservation of
admissibility, involutivity and compatibility:
\begin{equation}
 \begin{aligned}
 \rho_{n,m}C_n^{\mathrm b}&=C_m^{\mathrm b}\rho_{n,m},
 &t_m(C_m^{\mathrm b}x_m)&=C_Mt_m(x_m),\\
 \bar p_{n,m}\widetilde C_n&=\widetilde C_mp_{n,m},
 &\widetilde C_{m,\bar x}\widetilde C_{m,x}&=I.
 \end{aligned}
 \label{vi:vi:eq:levantamiento-cm}
\end{equation}
In the first row, the reader publishes the orientation register
of the history. The remaining transport of the state accompanies that
publication. The second row gives the lift laws;
the subscripts $x,\bar x$ make its two source fibres explicit.
If $s$ is persistent, its antisection is
$\bar s=(\widetilde C_ms_m)_m$.

\begin{proposition}[Persistence of the antisection]
\label{vi:vi:prop:antiseccion-persistente}
Conditions \eqref{vi:vi:eq:levantamiento-cm} transport persistent
sections to persistent sections and restore the section upon
conjugating twice. If they also transport the internal symmetries by
$g_m\mapsto\widetilde C_mg_m\widetilde C_m^{-1}$, they induce an
involution on the material classes of
Definition~\ref{vi:vi:def:particula}.
\end{proposition}
\begin{proof}
Compatibility gives
$\bar p_{n,m}\bar s_n=\widetilde C_mp_{n,m}s_n
=\widetilde C_ms_m=\bar s_m$.
Involutivity restores $s$. For $s'_m=g_ms_m$, their
antisections are related through
$\widetilde C_mg_m\widetilde C_m^{-1}$. The restriction law
in \eqref{vi:vi:eq:levantamiento-cm} preserves compatibility of this
family, which belongs to the internal groups by the condition of the
statement. The construction thus descends to the classes.
\end{proof}

The charge representation uses the odd part of the material
record. If $\mathfrak q$ is the stated charge character, its
law is $\mathfrak q(\bar s)=-\mathfrak q(s)$. The elementary charge
entering the expression of that coordinate on a dimensional line is
the internal output already constructed; the representation is subsequently applied
to the state without using a target charge to select its record.

For mass, covariance of its operator is required. Let
$\widehat M_s$ and $\widehat M_{\bar s}$ be the self-adjoint operators
of the two realisations, and let $\mathcal C$ be their unitary
or antiunitary intertwiner. Its domain is transported by
$\mathcal C\operatorname{Dom}(\widehat M_s)
=\operatorname{Dom}(\widehat M_{\bar s})$. The operator condition is
\begin{equation}
 \mathcal C\widehat M_s\mathcal C^{-1}
 =\widehat M_{\bar s}.
 \label{vi:vi:eq:covariancia-masa}
\end{equation}

\begin{proposition}[Spectral preservation under conjugation]
\label{vi:vi:prop:masa-conjugada}
Under \eqref{vi:vi:eq:covariancia-masa}, transport of an eigensection
with mass $m\in\mathbb R$ preserves that eigenvalue. The spectral
measure is transported by
$E_{\bar s}(B)=\mathcal C E_s(B)\mathcal C^{-1}$ for every real
Borel set $B$. In particular, a scalar realisation preserves
mass on passage to the antisection.
\end{proposition}
\begin{proof}
If $\widehat M_s\psi=m\psi$, covariance gives
$\widehat M_{\bar s}\mathcal C\psi
=\mathcal C(m\psi)=m\mathcal C\psi$; the last equality
also holds in the antiunitary case because $m$ is real. The real
functional calculus of self-adjoint operators transports their
spectral projectors and proves the second assertion.
\end{proof}

In a band record the same condition can be seen in elementary
form. Given a real reading $E_0(\tau)$, let
$E_\pm=(E_0\pm E_0\circ C_M)/2$ and let $\xi\in\{+1,-1\}$
be the sheet orientation. Then
\begin{equation}
 E_{\mathrm{band}}(\tau,\xi)=E_+(\tau)+\xi E_-(\tau),\qquad
 E_{\mathrm{band}}(C_M\tau,-\xi)=E_{\mathrm{band}}(\tau,\xi).
 \label{vi:vi:eq:banda-paridad}
\end{equation}
The equality follows from the evenness of $E_+$ and oddness of $E_-$.
It explains how band energy can be preserved by transporting
record and sheet simultaneously; its identification with the mass of a
particular realisation uses \eqref{vi:vi:eq:covariancia-masa}.

The electronic record
$\tau_e=(+,+,+,-,-,-,+,+,+,-,-,-)$ satisfies $C_M\tau_e=\tau_e$.
This equality concerns the twelve-position word. The pair
$(\tau_e,\xi)$ preserves an additional orientation changed by
conjugation. The distinction allows the unique electronic
centre and the conjugate charge realisations to be retained together;
a fixed word does not become a proof of neutrality.

\subsection{Sign holonomy and the persistence band}
\label{vi:vi:subsec:mobius}

Let $\ell$ be a return loop of the material reader of a surviving
extension. Its circuit group is $\langle\ell\rangle\simeq
\mathbb Z$. On this group, the sign character is specified by
\begin{equation}
 \varepsilon:\langle\ell\rangle\longrightarrow\{+1,-1\},
 \qquad \varepsilon(\ell)=-1.
 \label{vi:vi:eq:holonomia-signo}
\end{equation}
The circuit group records how many times the loop is traversed; its
type differs from the finite phase $C_9$. The interval bundle
associated with the character is
\begin{equation}
 \mathcal M_\ell=([0,1]\times[-1,1])/
                  \bigl((0,u)\sim(1,-u)\bigr).
 \label{vi:vi:eq:banda-mobius}
\end{equation}

\begin{theorem}[Oriented return on the band]
\label{vi:vi:thm:retorno-mobius}
The bundle \eqref{vi:vi:eq:banda-mobius} is a Möbius band.
Its transport over $\ell$ reverses the sign of the fibre and over
$\ell^2$ restores it. Thus a nonzero vector of that line
requires two circuits for its oriented return. This preserves the
distinction between return of orientation and return of the complete state.
\end{theorem}
\begin{proof}
The endpoint relation in \eqref{vi:vi:eq:banda-mobius} is the linear
map $u\mapsto-u$, with negative determinant. It gives the nonorientable
real line over the circle and its corresponding interval
band. Composing two transports produces $u\mapsto u$.
The memory of the extension is not part of this sign quotient;
its updating continues to obey TPK transport and
\eqref{vi:vi:eq:persistencia-fase-memoria}.
\end{proof}

Parametrising one circuit of the base by $2\pi$
represents restoration of the cover by $4\pi$. If the
temporal reader assigns $108$ steps to the relevant circuit, the
oriented lift closes at $216$ steps. These assignments
require the stated loop and reader; they do not divide the history into
two successive particles of $108$ steps.

\begin{proposition}[Decomposition of observables on the double cover]
\label{vi:vi:prop:observables-paridad}
Let $p:\widetilde X\to X$ be a principal double cover, with
involution $\iota$. For $f\in C^0(\widetilde X,\mathbb R)$,
\[
 f^+=\frac{f+f\circ\iota}{2},\qquad
 f^-=\frac{f-f\circ\iota}{2}
\]
give $f=f^++f^-$, with $f^+\circ\iota=f^+$ and
$f^-\circ\iota=-f^-$. The first part descends to $X$; the
second represents a section of the associated sign line. On
the Möbius line, every continuous global section has a zero.
\end{proposition}
\begin{proof}
The identities follow by substituting $\iota^2=I$. The invariant
part is constant on each fibre and defines a function on the
quotient. The anti-invariant part satisfies the transition law
of the sign representation. On the circle, a section of this
line is represented by a continuous real function $u$ on $[0,1]$
with $u(1)=-u(0)$. If an endpoint is zero, a zero already exists; otherwise
the endpoints have opposite signs, and the intermediate value
theorem provides an interior zero.
\end{proof}

The decomposition distinguishes two classes of observables on the
same cover. It represents odd orientation information
and preserved even data. Applying it to a material record uses
the lift \eqref{vi:vi:eq:levantamiento-cm}; the sign character
describes the transported line, not all the data of the
particle.

\subsection{Bidirectional lattice and channel conjugation}
\label{vi:vi:subsec:reticulo-bidireccional}

The integer angular register publishes two coordinates
\begin{equation}
 W_+=6n_A+s_C,\qquad W_-=6n_A-s_C,
 \qquad (n_A,s_C)\in\mathbb Z^2.
 \label{vi:vi:eq:reticulo-ida-retorno}
\end{equation}
The coefficients $n_A,s_C$ are internal route readings. The factor
$6$ belongs to the dodecaphase normalisation of that record.

\begin{proposition}[Lattice index and interchange of readings]
\label{vi:vi:prop:indice-bidireccional}
The image of \eqref{vi:vi:eq:reticulo-ida-retorno} has index $12$
in $\mathbb Z^2$ and Smith normal form $\operatorname{diag}(1,12)$.
It is characterised by $W_++W_-\equiv0\pmod{12}$. The operation
$s_C\mapsto-s_C$ interchanges $W_+$ and $W_-$.
\end{proposition}
\begin{proof}
The matrix of the map is
$\bigl(\begin{smallmatrix}6&1\\6&-1\end{smallmatrix}\bigr)$.
The greatest common divisor of its entries is $1$, and its determinant
is $-12$, giving the stated invariant factors. The inverse
coordinates are
$n_A=(W_++W_-)/12$ and $s_C=(W_+-W_-)/2$.
If the sum is divisible by $12$, the two integers have the same
parity and both expressions are integral. The condition is also
necessary by summing \eqref{vi:vi:eq:reticulo-ida-retorno}.
The interchange assertion is checked by substituting $-s_C$.
\end{proof}

In the real angular representation, the coordinates already generated,
$a=A_{\mathrm{rad}}$ and $b=C^*_{\mathrm{rad}}$, produce the
channels $q_\pm=\exp[-(a\pm b)]$. In the domain $a>|b|$, both
belong to $(0,1)$. Let $R^*=R$, $R^2=I$ be the sheet
grading and $C$ a unitary interchange with $CRC^{-1}=-R$.
For $T(b)=\exp[-(aI+bR)]$ we obtain
\begin{equation}
 CT(b)C^{-1}=T(-b),\qquad
 CP_\pm C^{-1}=P_\mp,\quad P_\pm=(I\pm R)/2.
 \label{vi:vi:eq:conjugacion-canales}
\end{equation}
Indeed, conjugation transforms the exponent and each term
of its convergent series. This proof transmits the interchange to
the channels. The integer lattice and the real operator have distinct
domains and publish the same interchange operation through
their declared reading maps.

\subsection{Spinor representation of the electronic block}
\label{vi:vi:subsec:spin-electron}

The three-dimensional APP construction of the addition and
multiplication projectors selects a principal electronic plane through the
tritic mode $(1,-1,0)^3$. Its derivation and the evaluation of the
projectors are included among the antecedents of the present
article. In an orthonormal basis of that plane, their restrictions
are
\begin{equation}
 P=\begin{pmatrix}1&0\\0&0\end{pmatrix},\qquad
 Q=\begin{pmatrix}1/4&\sqrt3/4\\\sqrt3/4&3/4\end{pmatrix},
 \qquad S=2P-I,\quad J=\frac4{\sqrt3}[P,Q].
 \label{vi:vi:eq:bloque-electronico}
\end{equation}
This representation receives the two constructed projectors and
preserves their orientation. Multiplication gives
$S^*=S$, $J^*=-J$, $S^2=I$, $J^2=-I$, and $SJ=-JS$.
Specifically, $S=\operatorname{diag}(1,-1)$ and
$J=\bigl(\begin{smallmatrix}0&1\\-1&0\end{smallmatrix}\bigr)$.

Let $\mathscr S_e=E_e\otimes_{\mathbb R}\mathbb C$ be the
complexification of the real plane. The scalar unit $i$ belongs
to the field of this representation; $J$ remains the real
endomorphism arising from the projectors.

\begin{proposition}[Hermitian triple and representation of rotations]
\label{vi:vi:prop:spin}
In $\mathscr S_e$, define
\begin{equation}
 \sigma_1=SJ,\qquad \sigma_2=-iJ,\qquad \sigma_3=S.
 \label{vi:vi:eq:terna}
\end{equation}
These operators are Hermitian and traceless, and satisfy
\begin{equation}
 \sigma_a\sigma_b=\delta_{ab}I+
 i\sum_c\epsilon_{abc}\sigma_c,
 \qquad L_a=\sigma_a/2,\qquad
 [L_a,L_b]=i\sum_c\epsilon_{abc}L_c,
 \quad\sum_aL_a^2=\tfrac34I.
 \label{vi:vi:eq:spin-algebra}
\end{equation}
They constitute the irreducible spin-$1/2$ representation on
the electronic block.
\end{proposition}
\begin{proof}
The relations of $S,J$ give $(SJ)^*=SJ$, $(-iJ)^*=-iJ$, and
$\sigma_a^2=I$. The cyclic products are
$(SJ)(-iJ)=iS$, $(-iJ)S=iSJ$, and $S(SJ)=J=i(-iJ)$;
reversing the order reverses the sign. This proves the product
identity and the commutators. In the basis of
\eqref{vi:vi:eq:bloque-electronico}, the three matrices are
\[
 \begin{pmatrix}0&1\\1&0\end{pmatrix},\quad
 \begin{pmatrix}0&-i\\i&0\end{pmatrix},\quad
 \begin{pmatrix}1&0\\0&-1\end{pmatrix}.
\]
Their trace is zero and, together with $I$, they generate $M_2(\mathbb C)$.
A subspace invariant under all three is invariant under that entire
algebra and can only be $0$ or $\mathscr S_e$. Finally,
summing $L_a^2=I/4$ gives the value $3I/4$ of the quadratic operator.
\end{proof}

Recognition of this triple as the Pauli matrices follows
\eqref{vi:vi:eq:bloque-electronico}. The norm and product
$\langle X,Y\rangle=\tfrac12\operatorname{tr}(XY)$ on the
real space of traceless Hermitian matrices give its space
of directions. For $n\in\mathbb R^3$, $|n|=1$, let
\begin{equation}
 H_e(n)=\sum_an_a\sigma_a,\qquad
 h_e(n)=\tfrac12H_e(n),\qquad
 \Pi_\pm(n)=\tfrac12(I\pm H_e(n)).
 \label{vi:vi:eq:helicidad}
\end{equation}

\begin{proposition}[Directional decomposition and double return]
\label{vi:vi:prop:helicidad}
The operators $\Pi_\pm(n)$ are rank-one orthogonal
projectors and $h_e(n)\Pi_\pm(n)=\pm\tfrac12\Pi_\pm(n)$.
The rotation group around $n$ is
\begin{equation}
 U_n(\theta)=e^{-i\theta H_e(n)/2}
 =\cos(\theta/2)I-i\sin(\theta/2)H_e(n),\qquad
 U_n(2\pi)=-I,\quad U_n(4\pi)=I.
 \label{vi:vi:eq:retorno-spin}
\end{equation}
\end{proposition}
\begin{proof}
The antisymmetric part of \eqref{vi:vi:eq:spin-algebra} vanishes
when contracted with $n_an_b$, so $H_e(n)^2=I$.
Idempotence, orthogonality, and the sum
$\Pi_++\Pi_-=I$ follow. They are Hermitian with trace one, hence their
rank is one. Multiplication by $H_e(n)/2$ gives the eigenvalues.
Separating even and odd powers of the exponential produces
\eqref{vi:vi:eq:retorno-spin} and its group law. Angular normalisation
is checked by
\[
 U_n(\theta)H_e(v)U_n(\theta)^*
 =H_e\bigl(v\cos\theta+(n\times v)\sin\theta
       +n(n\cdot v)(1-\cos\theta)\bigr),
\]
which follows from multiplication using the product identity of the
triple. The induced action on directions is a rotation through
angle $\theta$; the spinor lift carries its half-angle.
\end{proof}

If $n$ is realised as the direction of propagation, $h_e(n)$ is
dimensionless helicity. Reversal of that direction satisfies
$H_e(-n)=-H_e(n)$ and $\Pi_\pm(-n)=\Pi_\mp(n)$.
This operation changes the directional labels; charge conjugation
transports the material record through
\eqref{vi:vi:eq:levantamiento-cm}. Chirality is defined through
the relativistic grading appropriate to the representation and
its projectors, which differ from \eqref{vi:vi:eq:helicidad}.

The scalar electronic evaluation acts as
$\varepsilon_e^\beta I$ on $\mathscr S_e$ and commutes with
the triple and with $\Pi_\pm(n)$. Rotational information preserves
the structure of that evaluation without introducing another factor into
mass. The Möbius band, in turn, transports a real
line around a loop, and nonadic holonomy transports
phase and memory. The three returns can enter the same
realisation; their domains and intertwiners determine the
precise composition.
